\documentclass[twoside]{article}

\usepackage{morewrites}

\usepackage[svgnames,dvipsnames,x11names]{xcolor}
\usepackage{graphicx}
\usepackage[english]{babel}
\usepackage[colorlinks=true, linkcolor=NavyBlue, urlcolor=NavyBlue, citecolor=NavyBlue]{hyperref}
\usepackage[margin=15mm]{geometry}
\usepackage{ts-fonts}
\usepackage{tcolorbox}
\usepackage{xcolor}
\usepackage{float}
\usepackage{seqsplit}
\usepackage{ts-callouts}

\usepackage{lastpage}
\usepackage{refcount}
\makeatletter
\def\ps@plain{
\let\@oddhead\@empty
\let\@evenhead\@empty
\def\@oddfoot{
\hfil
\ifnum\getpagerefnumber{LastPage}>1\relax
\thepage
\fi
\hfil}
\let\@evenfoot\@oddfoot
}
\makeatother

\title{Admonitions / Callouts\\[0.5em]\large An Overview of \textbf{classic} framed elements}
\author{}\date{}
\setlength{\emergencystretch}{3em}
\widowpenalties 3 10000 10000 150
\clubpenalties  3 10000 10000 150
\displaywidowpenalty 10000
\renewenvironment{abstract}{
\small
\begin{center}{\bfseries\abstractname}\end{center}
\begin{quotation}
}{\end{quotation}}
\makeatletter
\newcount\ts@affilcnt
\newcount\ts@loopcnt
\renewcommand{\thanks}[1]{
\stepcounter{footnote}
\protected@xdef\@thefnmark{\thefootnote}
\@footnotemark
\global\advance\ts@affilcnt\@ne
\expandafter\xdef\csname ts@affilM\the\ts@affilcnt\endcsname{\@thefnmark}
\expandafter\gdef\csname ts@affilT\the\ts@affilcnt\endcsname{#1}
}
\newcommand\tsreplayfn{
\ts@loopcnt\z@
\loop
\advance\ts@loopcnt\@ne
\ifnum\ts@loopcnt>\ts@affilcnt\else
\xdef\@thefnmark{\csname ts@affilM\the\ts@loopcnt\endcsname}
\@footnotetext{\csname ts@affilT\the\ts@loopcnt\endcsname}
\repeat
}
\makeatother
\begin{document}
\twocolumn[{
\maketitle
}]
\tsreplayfn
\begin{tscallout}[kind=note]
Highlights extra information that's useful but not critical, helping readers understand nuances.
\end{tscallout}
\begin{tscallout}[kind=tip]
Offers a practical hint that makes the task easier, saving the reader time or effort.
\end{tscallout}
\begin{tscallout}[kind=warning]
Draws attention to something that could cause problems if ignored, helping avoid mistakes.
\end{tscallout}
\begin{tscallout}[kind=caution]
Signals a potentially risky action, encouraging the reader to proceed carefully.
\end{tscallout}
\begin{tscallout}[kind=important]
Emphasizes key information the reader must not overlook to ensure proper understanding.
\end{tscallout}
\begin{tscallout}[kind=danger]
Flags a serious hazard that could break things or cause real harm if mishandled.
\end{tscallout}
\begin{tscallout}[kind=info]
Provides neutral, factual context that supports understanding without urgency.
\end{tscallout}
\begin{tscallout}[kind=hint]
Gives a gentle clue that helps the reader figure something out without revealing everything.
\end{tscallout}
\begin{tscallout}[kind=seealso]
Points to related material so the reader can explore deeper or connected topics.
\end{tscallout}
\begin{tscallout}[kind=question]
Raises an inquiry that prompts reflection or introduces a point the reader should consider.
\end{tscallout}
\begin{tscallout}[kind=abstract]
Summarizes the core ideas to help the reader grasp the purpose of a section or document quickly.
\end{tscallout}
\begin{tscallout}[kind=summary]
Recaps key points so the reader can retain the most important information at a glance.
\end{tscallout}
\begin{tscallout}[kind=success]
Celebrates a positive outcome or achievement, reinforcing good practices and results.
\end{tscallout}
\begin{tscallout}[kind=failure]
Highlights a setback or error, helping the reader learn from mistakes and avoid them in the future
\end{tscallout}
\begin{tscallout}[kind=bug]
Identifies a known issue or problem, guiding the reader on what to watch out for.
\end{tscallout}
\begin{tscallout}[kind=quote]
Presents a relevant quotation that adds authority or perspective to the content.
\end{tscallout}
\section*{Custom Admonition}
\begin{tscallout}[kind=unicorn]
A custom admonition with a unicorn theme, adding a whimsical touch to the information presented.
\end{tscallout}
\end{document}
